% Part: first-order-logic
% Chapter: proof-systems
% Section: axiomatic-deduction

\documentclass[../../../include/open-logic-section]{subfiles}

\begin{document}

\iftag{FOL}
      {\olfileid{fol}{prf}{axd}}
      {\olfileid{pl}{prf}{axd}}

\olsection{مسلّمی \usetoken{P}{derivation}}

مسلّمی !!{derivation}s سب توں پرانے تے سب توں سادہ منطقی
!!{derivation} نظام نیں۔ ایہناں دے !!{derivation}s بس !!{sentence}s دے
سلسلے ہُندے نیں۔ !!{sentence}s دا کوئی سلسلہ اک درست !!{derivation}
اُدوں منّیا جاندا اے جدوں اوہدے وچ ہر !!{sentence}~$!A$ ہِٹھ لکھیاں
شرطاں وچوں کوئی اک پوری کرے:
\begin{enumerate}
\item $!A$ اک مسلّمہ اے، یا
\item $!A$، !!{sentence}s دے دتے ہوئے سیٹ~$\Gamma$ دا !!a{element} اے، یا
\item $!A$ نوں استنباط دے کسے قاعدے توں جواز ملدا اے۔
\end{enumerate}
مسلّمہ ہون لئی $!A$ نوں کجھ مقررہ !!{sentence} خاکیاں وچوں کسے اک دا
روپ رکھنا پیندا اے۔ مسلّمہ خاکیاں دے بہت سارے ایہو جہے سیٹ نیں جیہڑے
پہلے درجے دی منطق لئی اک تسلی بخش (صحت مند تے مکمل) !!{derivation} نظام
دیندے نیں۔ کجھ نوں اوہناں رابطیاں دے لحاظ نال ترتیب دتا جاندا اے
جنہاں اُتے اوہ لاگو ہُندے نیں؛ مثال لئی، خاکے
\[
!A \lif (!B \lif !A) \qquad !B \lif (!B \lor !C) \qquad (!B \land !C) \lif !B
\]
عام مسلّمے نیں جیہڑے $\lif$، $\lor$ تے~$\land$ اُتے لاگو ہُندے نیں۔
کجھ مسلّمی نظاماں دا مقصد مسلّماں دی گنتی گھٹ توں گھٹ رکھنا ہُندا
اے۔ ایہ گل اوس اُتے منحصر اے کہ کیہڑے رابطیاں نوں ابتدائی منّیا
جاندا اے؛ ایسے مسلّمی نظام لبھنے وی ممکن نیں جیہڑے صرف اک مسلّمے
اُتے مشتمل ہون۔

استنباط دا قاعدہ اک شرطیہ عبارت اے جیہڑی !!a{derivation} وچ
!!a{sentence} نوں جواز ملن لئی کافی شرط دیندی اے۔ قاعدۂ اثباتِ مقدم
ایہو جہا اک بہت عام قاعدہ اے: ایہ آکھدا اے کہ جے $!A$ تے $!A \lif
!B$ نوں پہلاں ای جواز مل چکیا اے، تاں $!B$ نوں وی جواز ملدا اے۔ ایس
دا مطلب اے کہ !!a{derivation} وچ !!{sentence}~$!B$ والی سطر نوں جواز
ملدا اے، بشرطیکہ $!A$ تے $!A \lif !B$ دونیں (کسے
!!{sentence}~$!A$ لئی) اوس !!{derivation} وچ~$!B$ توں پہلاں آ چکے ہون۔

مسلّمی !!{derivation}s اُتے مبنی $\Proves$ تعلق دی تعریف ایویں کیتی
جاندی اے: $\Gamma \Proves !A$ اُدوں تے صرف اُدوں جدوں کوئی
!!a{derivation} ایہو جہا ہووے کہ !!{sentence}~$!A$ اوہدا آخری فارمولا
ہووے (تے $\Gamma$ نوں اوس !!{derivation} دے اوہناں !!{sentence}s دا
سیٹ منّیا جاوے جنہاں نوں اُتّے دتی شرط~(2) توں جواز ملدا اے)۔ $!A$
اک قضیہ اے جے $!A$ دا کوئی ایہو جہا !!a{derivation} ہووے جتھے
$\Gamma$ خالی ہووے، یعنی !!{derivation} وچ ہر !!{sentence} نوں یا
(1) توں یا~(3) توں جواز ملدا ہووے۔ مثال لئی، ایہ !!a{derivation}
وکھاؤندا اے کہ $\Proves !A  \lif (!B \lif (!B \lor !A))$:
\begin{derivation}
  1. & $!B \lif (!B \lor !A)$ \\
  2. & $(!B \lif (!B \lor !A)) \lif (!A  \lif (!B \lif (!B \lor !A)))$\\
  3. & $!A  \lif (!B \lif (!B \lor !A))$
\end{derivation}
سطر~1 دا !!{sentence} مسلّمے $!A \lif (!A \lor !B)$ دے روپ دا اے
($!A$ تے $!B$ دے کردار الٹے کر کے)۔ سطر~2 دا جملہ مسلّمے $!A \lif
(!B \lif !A)$ دے روپ دا اے۔ ایس لئی دوناں سطراں نوں جواز ملدا اے۔
سطر~3 نوں قاعدۂ اثباتِ مقدم توں جواز ملدا اے: جے اسیں ایہنوں مختصر
کر کے $!D$ لکھیے، تاں سطر~2 دا روپ $!C \lif !D$ اے، جتھے $!C$،
$!B \lif (!B \lor !A)$ یعنی سطر~1 اے۔

سیٹ $\Gamma$ متضاد اے جے $\Gamma \Proves \lfalse$۔ اک مکمل مسلّمی
نظام $\lfalse \lif !A$ نوں وی، ہر~$!A$ لئی، ثابت کرے گا، ایس لئی جے
$\Gamma$ متضاد اے تاں $\Gamma \Proves !A$، ہر~$!A$ لئی۔

منطق لئی مسلّمی !!{derivation}s نظام سب توں پہلاں Gottlob Frege نے
اپنی 1879 دی \emph{Begriffsschrift} وچ دتے سن؛ ایس کارن اوہ کم اکثر
جدید منطق دا پہلا کم منّیا جاندا اے۔ Alfred North Whitehead تے
Bertrand Russell دی \emph{Principia Mathematica} وچ، تے 1920 دی دہائی
وچ David Hilbert تے اوہدے شاگرداں ولوں، ایہناں نظاماں نوں کمال تک
پہنچایا گیا۔ ایس لئی ایہناں نوں اکثر ``فریگ نظام'' یا ``ہلبرٹ نظام''
آکھیا جاندا اے۔ ایہ بہت لچکدار نیں، ایس معنی وچ کہ کسے منطق لئی
مسلّمی نظام لبھنا اکثر سَوکھا ہُندا اے۔ چونکہ !!{derivation}s دی بناوٹ
بہت سادہ اے تے اوہناں وچ استنباط دے صرف اک یا دو قاعدے ہُندے نیں، ایس
لئی اوہناں دے \emph{بارے} گلاں ثابت کرنا وی نسبتاً سَوکھا اے۔ پر
عملی ورتوں وچ ایہ بہت اوکھے نیں، یعنی ثبوت لبھنے تے لکھنے اوکھے نیں۔

\end{document}
